\chapter{From molecular computation to genomic computation}
\label{ch:dna32}
\section{Finish the conceptual arc without erasing boundaries}
A DNA molecule can be a physical reactant, a sequence record, a formal string
or an encoded tensor. These descriptions are useful precisely because they
are not identical. This closing chapter establishes what can be transferred
between them, what cannot be reconstructed after projection, and what would
make a genomic-computation hypothesis scientifically informative.

Chapter 31 connected specification, realization, observation and scoped evidence.
We now apply that discipline to the boundary between molecular mechanisms and
software models inspired by genomic organization. The original executable
examples are constructive information-loss witnesses, strict encodings and
paired evaluation arithmetic. No molecular experiment or trained-model result
is reported here.

This is the final planned chapter of the 32-chapter DNA authoring spine.
It is not a claim that the entire publication has passed independent review.
Evolutor continues beyond this boundary with testable architectural and systems
hypotheses.

\section{A projection keeps a record, not the whole molecule}
Write the representation chain
\[
 M\xrightarrow{\pi_1}S\xrightarrow{\pi_2}\Sigma^*
 \xrightarrow{\phi}X\xrightarrow{F}Y.
\]
$M$ contains physical molecular state; $S$ is a sequence record with coordinate
conventions; $\Sigma^*$ is a formal word; $X$ is a numerical encoding; $F$ is a
software operator. Each arrow has its own domain, codomain and convention.
A bare sequence omits concentration, copy count, solvent, modifications,
geometry and interaction partners. An orientation frame may be carried as
metadata or lost when only base order is retained.
\Plate{01-ladder}{Representations and their information contracts. Encoding a physical description is not a physical execution of the encoded algorithm.}{fig:dna32-ladder}

\section{Prove when a discarded property cannot be recovered}
Let $\pi:M\to R$ be a projection and $q:M\to Q$ a property. A function
$\widehat q:R\to Q$ satisfying $\widehat q(\pi(m))=q(m)$ for every admissible
$m$ exists exactly when $q$ is constant on every fiber of $\pi$.

For necessity, if $\pi(m_1)=\pi(m_2)$, applying $\widehat q$ gives
$q(m_1)=q(m_2)$. For sufficiency, assign to each reachable representation
$r$ the common value of $q$ among objects mapping to it. The fiber condition
makes that assignment well-defined; unreachable representations can be handled
outside the admissible domain.

This is a \Term{fiber recoverability criterion}. It is stronger than a vague
claim that information is lost: it says precisely which property cannot be
determined from a chosen record.
\Plate{02-ledger}{A constructive collision: identical base order and frame, different concentration. A sequence-only estimator must return one answer for both and cannot be correct for both.}{fig:dna32-ledger}
\Listing{collision_witness}

The code searches a finite list for such a witness. Finding one disproves
recoverability on that domain. Failing to find one does not prove universal
injectivity; the tested samples may simply omit a collision.

\subsection{A kinetic consequence, not only a metadata consequence}
Take samples with the same base order and orientation but $F=1$ and $F=10$
in assigned concentration units. With $k_{\rm on}=.1$ and $A=.2$, the capture
rates $k_{\rm on}AF$ are $.02$ and $.2$. No algorithm reading only the identical
sequence records can infer both correct rates.

An external calibrated concentration measurement could resolve the distinction.
That is new evidence with provenance and context, not recovery of information
that the sequence already contained. Chapter 31's observation contract applies
to that measurement.

\section{Use a preservation ledger with the right scope}
For a representation, record the named properties it declares to carry.
A pure projection may discard a property but must not silently introduce one.
\Listing{project}
The property-set checker is a \Term{coverage ledger}, not a truth verifier.
It cannot certify that a file actually uses the declared coordinate convention
or that a supplied concentration was measured correctly.

Derived properties also require care. Reverse complement can be computed from
a canonical string by an explicit algorithm; its addition is a derived result,
not evidence that the original string encoded concentration. New observations,
derivations and lossy projections should have distinct provenance records.

\section{An exact encoding is exact only at its declared level}
For alphabet order $(A,C,G,T)$, one-hot maps each base to a basis vector in
$\mathbb R^4$. A strict integer one-hot array preserves canonical base order:
\[
 \operatorname{decode}(\operatorname{encode}(s))=s.
\]
Each row must have four integer entries, exactly one equal to one.
The reference rejects Boolean and floating-point lookalikes rather than
silently weakening its data contract.
\Listing{decode_one_hot}

A naked one-hot tensor does not preserve an external orientation frame unless
that frame is separately carried. Nor does it preserve the molecule's geometry.
The reverse-complement identity $\operatorname{rc}(\operatorname{rc}(s))=s$
is an exact string statement. It does not describe a molecule spontaneously
executing a reversal operation. Exhaustive tests over canonical strings through
length five independently construct complements and verify round trips.

Learned embeddings can intentionally collapse distinct bases or contexts.
Compression may aid a task while losing exact recoverability. The correct
question is whether the retained representation is sufficient for the specified
task, not whether all information must always be kept.

\section{Sufficiency is task-dependent}
For a task target $q$, the fiber criterion tests deterministic sufficiency:
$q$ must factor through the representation. A probabilistic prediction instead
requires a distribution over omitted states and its domain of validity.
A learned conditional average is not an exact reconstruction of each underlying
molecule.

For example, two concentrations might have the same classification label but
different kinetic trajectories. The sequence could be sufficient for the former
under a declared labeling rule and insufficient for the latter. A claim of
sufficiency must name the task and domain.

\section{Engineered reactions and genomic regulation}
An engineered circuit selects reactants, interfaces and resource assumptions.
A living genomic system involves cis elements, transcription factors, chromatin,
spatial contacts, RNA/protein turnover, signaling and developmental history.
Both involve nucleic acids, but their causal models differ.
\Plate{03-two-worlds}{Contrasting mechanistic organization. Designed displacement and genomic regulation share a substrate vocabulary, not an automatically shared computation.}{fig:dna32-worlds}

A DNA sequence can influence a reaction without being the entire reactive state.
Likewise, a promoter sequence contributes to regulation without determining
expression independently of cellular context. Chapter 30's regulatory and
developmental distinctions must not be collapsed into a single software gate.
Earlier molecular work, including Adleman's experiment and formal-to-DNA
kinetic constructions, establishes specific physical/programming relationships,
not universal equivalence between molecules and digital models
\cite{dna32-adleman,dna32-soloveichik}.

\section{Turn inspiration into an operator and prediction}
A useful transfer names a biological mechanism, identifies a computational
property, formalizes an operator, compares it with the closest established
mechanism and proposes a falsifiable intervention.

A software sigmoid is already a standard gate. Calling it a promoter does not
add a mechanism. A proposal must specify an actual difference: typed multilevel
control, strand symmetry, structural generation, locality, multiple timescales
or another explicit property. Then it must explain why that difference should
alter a declared observable under matched conditions.

Development-inspired structural generation must name inherited specification,
transition rule, environment, clock, persistence and learning signal.
Population-based evolution requires population, heritable variation and
differential selection; gradient descent in one parameter vector is not
automatically evolutionary search.
\Plate{04-loop}{A matched $2\times2$ intervention design. Mechanism removal and transplant provide different tests of the claimed explanatory role.}{fig:dna32-loop}

\section{Compute ablation and transplant effects}
Let $Y_{b,g,s}$ be a metric for backbone $b\in\{0,1\}$, mechanism
$g\in\{0,1\}$ and paired seed $s$. Define
\[
 d^{\rm transplant}_s=Y_{0,1,s}-Y_{0,0,s},\qquad
 d^{\rm ablation}_s=Y_{1,1,s}-Y_{1,0,s}.
\]
The interaction is
$d^{\rm int}_s=d^{\rm ablation}_s-d^{\rm transplant}_s$.
It distinguishes a mechanism's effect in the target backbone from its effect
when transplanted into the baseline backbone. A strong full model can coexist
with zero effect from the named mechanism.
\par\Needspace{16\baselineskip}
\Listing{paired_effect}
\Listing{factorial_effect}

The following scores are deliberately assigned synthetic arithmetic, not
training runs, benchmark results or evidence of any architecture:
\begin{center}\input{\Generated/panel.tex}\end{center}
They give mean transplant effect $.0275$, ablation effect $.0375$ and
interaction $.01$. Pair differences by seed before computing their spread.
For $n>1$, sample standard deviation is
$s_d=\sqrt{\sum_s(d_s-\bar d)^2/(n-1)}$ and the usual descriptive standard-error
calculation is $s_d/\sqrt n$. With one pair, those quantities are unidentified.
No confidence level is asserted by this computation.

\section{What makes this a causal experiment?}
Arithmetic alone does not identify causal effects. The mechanism interventions
must be controlled; data splits, optimizer settings, training budget, parameter
and activation budgets, evaluation protocol and seed pairing must be declared.
Ablation can alter capacity or optimization difficulty; transplant can require
different interfaces. These are potential confounds, not reasons to omit the
comparison.

Predeclare a primary metric, adverse outcomes, exclusion criteria and failure
thresholds. Separate tuning from held-out evaluation. Use matched controls
where possible, including a structure-preserving but semantically disrupted
mechanism. Report all runs and the metric direction: a positive difference in
loss is not an improvement merely because it is positive.

An effect in a synthetic motif task does not establish genomic understanding,
continual learning or general intelligence. Out-of-domain biology and software
capability each require their own tasks and observations.

\section{Keep the two model directions distinct}
\Plate{05-programs}{Program boundaries. DOGMA is the non-Transformer research direction; HermonDNA is the Transformer branch; Evolutor supplies the broader research and systems program.}{fig:dna32-programs}
DOGMA and HermonDNA must earn support through explicit operators and comparisons,
not DNA vocabulary. Evolutor owns the broader hypotheses, evaluation, compiler
and runtime contracts. DNA Computing supplies molecular, formal and evidential
foundations; it does not pre-decide which architectural proposal will succeed.

A causal information-access contract matters in either branch. Offline suffix
state cannot be used to claim causal decoding. A restricted affine scan cannot
establish equivalence to an arbitrary state-dependent architecture. These
boundaries remain part of the handoff.

\section{Replace the ladder with scoped evidence}
Chapter 31 rejected an ordinal ordering of formal, simulated, cell-free and
cellular evidence. The final chapter must retain that correction.
A kinetic theorem, simulated trajectory, fluorescence measurement and neural
benchmark are different evidence branches, each attached to a claim and context.

Transferring between branches requires an explicit bridge. A cell experiment
does not prove a software architecture; a neural benchmark does not validate a
molecular rate. Listing more advanced environments cannot substitute for
identifiability, controls or independent replication.

\section{Close the authoring spine, not the publication gates}
\Plate{06-map}{The book's closing map and separate release gates. Completed chapter authoring is not independent scientific acceptance.}{fig:dna32-map}
The manuscript now reaches the planned conceptual endpoint: physical resources
and mechanisms; formal languages and algorithms; strand displacement and
chemical networks; reliability and readout; symbolic, stochastic and
differentiable models; genomic organization; evidence and abstraction contracts.

The active continuous edition must still receive cumulative editorial review:
notation, bibliography, glossary, index, cross-chapter assumptions and visual
consistency. Independent scientific and reader review remain distinct gates.
Legacy accepted material stays preserved; the word complete refers here to
the planned chapter spine, not a certified publication.

\section{Twenty worked synthesis exercises}
\begin{enumerate}
\item \textbf{Projection.} Two molecules have one sequence but different
concentrations. \emph{Answer:} sequence projection is many-to-one.
\item \textbf{Recovery.} State the fiber criterion. \emph{Answer:}
the desired property must be constant over every projection fiber.
\item \textbf{Proof.} Why does one collision disprove exact recovery?
\emph{Answer:} the estimator receives one representation but would need two
different outputs.
\item \textbf{Finite search.} No collision was found in ten samples.
\emph{Answer:} this is not proof of universal recoverability.
\item \textbf{Kinetics.} Compare $kAF$ for $k=.1,A=.2,F=1$ versus $10$.
\emph{Answer:} $.02$ versus $.2$.
\item \textbf{New measurement.} Can calibrated concentration metadata help?
\emph{Answer:} yes, as newly sourced evidence, not recovery from sequence alone.
\item \textbf{One-hot.} What alphabet order is used? \emph{Answer:}
$(A,C,G,T)$, explicitly fixed.
\item \textbf{Invalid row.} Is $(1,1,0,0)$ legal? \emph{Answer:} no;
exactly one integer one is required.
\item \textbf{Frame.} Does a naked tensor preserve an external orientation
label? \emph{Answer:} no; carry the label separately.
\item \textbf{Reverse complement.} Where is involution exact?
\emph{Answer:} canonical strings, not a claim of physical execution.
\item \textbf{Task sufficiency.} Can a representation suffice for one label
but not a trajectory? \emph{Answer:} yes; sufficiency is task-specific.
\item \textbf{Biological naming.} Is a sigmoid biologically a promoter?
\emph{Answer:} no; an operator and evidence bridge are required.
\item \textbf{Development.} What must structural generation specify?
\emph{Answer:} inherited rule, state, signals, clock, persistence and training.
\item \textbf{Evolution.} Is single-vector gradient descent population search?
\emph{Answer:} not without population, inheritance, variation and selection.
\item \textbf{Transplant.} Find the worked mean effect. \emph{Answer:} $.0275$.
\item \textbf{Ablation.} Find the target effect. \emph{Answer:} $.0375$.
\item \textbf{Interaction.} Subtract those effects. \emph{Answer:} $.01$.
\item \textbf{Causality.} Do assigned scores establish a mechanism's advantage?
\emph{Answer:} no; they demonstrate arithmetic only.
\item \textbf{Taxonomy.} Which branch is Transformer-based?
\emph{Answer:} HermonDNA, not DOGMA.
\item \textbf{Completion.} Does Chapter 32 confer publication acceptance?
\emph{Answer:} no; cumulative editorial, scientific and reader review remain.
\end{enumerate}

\section{Closing boundary}
Computation is a representation, a mechanism, a resource budget and an
observation contract, not only an equation. Genomic computation should inherit
that discipline before it inherits biological names. That is the book's
handoff to Evolutor.
